% This file was converted to LaTeX by Writer2LaTeX ver. 1.9.9
% see http://writer2latex.sourceforge.net for more info
\documentclass{article}
\usepackage{calc,amsmath,amssymb,amsfonts}
\usepackage[LGR,T1]{fontenc}
\usepackage[greek,italian]{babel}
\usepackage[style=numeric,backend=biber]{biblatex}
\author{HP}
\date{2026-07-26}
\begin{document}
\section{PROLOGO}
\subsection[Ci ho messo quarant{}'anni]{Ci ho messo quarant'anni}
Ci sono oggetti che, con il passare del tempo, smettono di essere semplici oggetti.

Diventano contenitori di ricordi.

Una fotografia ingiallita. Un vecchio quaderno di scuola. Una musicassetta con un'etichetta scritta a mano. Una scatola
dimenticata in soffitta.

La maggior parte delle persone, aprendola, vede soltanto plastica impolverata, cavi ingialliti e qualche vecchio
manuale. Per altri, invece, quella scatola è una macchina del tempo.

Dentro ce n'era una anche per me.

L'ho ritrovata molti anni dopo. Non ricordo nemmeno cosa stessi cercando quel giorno. Forse una vecchia valigia. Forse
degli album fotografici. Ricordo soltanto che, spostando alcuni scatoloni, è comparsa lei.

Era esattamente dove l'avevo lasciata.

Un po' più impolverata.

Un po' più silenziosa.

Ma era ancora lì.

L'ho aperta lentamente.

Per un istante ho avuto la sensazione che il tempo si fosse fermato.

Il primo oggetto che ho riconosciuto è stato il manuale. Aveva ancora quella copertina che conoscevo a memoria. Sotto
c'erano alcuni cavi, il Datasette, un joystick consumato dall'uso e, infine, lui.

Il Commodore 64.

L'ho preso tra le mani con la stessa attenzione con cui si prende in mano qualcosa di fragile. Non perché temessi di
romperlo. Ma perché, all'improvviso, mi sono reso conto che non stavo toccando soltanto un computer.

Stavo toccando una parte della mia infanzia.

Mi sono seduto.

Sono rimasto qualche minuto a guardarlo senza dire una parola.

Poi ho sorriso.

$\text{\textgreek{«}}$Ci ho messo quarant'anni.$\text{\textgreek{»}}$

Non quarant'anni per ritrovare il Commodore.

Quarant'anni per mantenere una promessa.

Una promessa fatta da un bambino che frequentava la quinta elementare e che, un giorno, aveva convinto i propri genitori
a fare un sacrificio economico enorme pur di portare a casa quella macchina color crema.

Quel bambino aveva promesso che quel computer gli avrebbe cambiato la vita.

Aveva ragione.

Gliel'ha cambiata davvero.

Grazie a quel Commodore ho scoperto il BASIC. Ho imparato a ragionare come un programmatore. Ho scelto gli studi che
avrei intrapreso. Ho trasformato una curiosità in una professione e oggi, dopo tanti anni, continuo ancora a fare il
mestiere che avevo immaginato allora.

Eppure c'era una promessa che non avevo mai mantenuto.

Il motivo per cui avevo desiderato quel computer non era imparare il BASIC.

Non era nemmeno giocare.

Io volevo capire cosa ci fosse dietro un videogioco.

Volevo scoprire come fosse possibile che un insieme di pixel riuscisse a raccontare un'avventura, a farmi emozionare, a
trasportarmi in mondi lontani.

Volevo crearne uno.

Non ci sono mai riuscito.

Per la verità, non l'ho mai nemmeno iniziato.

Con il mio amico Emanuele avevamo trascorso pomeriggi interi a immaginare il nostro primo videogioco. Gli avevamo
persino trovato un nome. Ne parlavamo come se fosse già pronto, ma ogni volta ci mancava qualcosa. Un'idea. Un
personaggio. Una storia. O forse semplicemente il coraggio di cominciare.

Poi la vita ha preso velocità.

Sono arrivati gli studi, il lavoro, gli impegni, le responsabilità.

Il Commodore è finito in quella scatola.

Il sogno è rimasto lì con lui.

Fino a oggi.

Questo libro nasce per un motivo molto semplice.

Voglio chiudere un cerchio rimasto aperto per oltre quarant'anni.

Ma non voglio farlo da solo.

Vorrei che, da questo momento in poi, tu ti sedessi accanto a me.

Non importa se hai ancora il tuo Commodore 64 perfettamente funzionante, se lo hai ritrovato in cantina dopo decenni o
se lo conoscerai soltanto attraverso un emulatore. Non importa nemmeno se negli anni Ottanta eri già un ragazzo curioso
o se sei nato molto tempo dopo.

Quello che conta è la voglia di scoprire.

Nel nostro laboratorio ci sarà anche un altro compagno di viaggio.

Si chiama ChatGPT.

Non scriverà il videogioco al nostro posto.

Non prenderà decisioni al posto nostro.

Farà quello che, negli anni Ottanta, faceva l'amico che sembrava sapere sempre qualcosa in più: suggerirà un'idea, farà
una domanda, ci aiuterà a osservare un problema da un punto di vista diverso.

Il resto toccherà a noi.

Perché il bello della programmazione non è ricevere una risposta.

È vivere quel momento meraviglioso in cui, dopo tanti tentativi, qualcosa finalmente funziona.

Alla fine di questo viaggio avremo costruito un videogioco per Commodore 64.

Ma, se tutto andrà come spero, quello sarà soltanto il risultato più visibile.

Il vero traguardo sarà un altro.

Riscoprire il bambino che, una sera del 1983, guardava una tastiera color crema e pensava che, da qualche parte, dietro
tutti quei tasti, fosse nascosto un piccolo segreto.

È arrivato il momento di scoprirlo.

Insieme.


\bigskip
\end{document}
